\documentclass[10pt,a4paper]{amsart}
\usepackage{amsmath,amssymb,amsthm,graphicx}
\usepackage{verbatim}
\newenvironment{reference}{\comment}{\endcomment}
%\newenvironment{reference}{}{}
\newenvironment{slogan}{\comment}{\endcomment}
\newenvironment{history}{\comment}{\endcomment}
% For commutative diagrams we use Xy-pic
\usepackage[all]{xy}

% We use 2cell for 2-commutative diagrams.
\xyoption{2cell}
% Theorem environments.
%
\theoremstyle{plain}
\newtheorem{theorem}[subsection]{Theorem}
\newtheorem{proposition}[subsection]{Proposition}
\newtheorem{lemma}[subsection]{Lemma}

\theoremstyle{definition}
\newtheorem{definition}[subsection]{Definition}
\newtheorem{example}[subsection]{Example}
\newtheorem{exercise}[subsection]{Exercise}
\newtheorem{situation}[subsection]{Situation}

\theoremstyle{remark}
\newtheorem{remark}[subsection]{Remark}
\newtheorem{remarks}[subsection]{Remarks}

\numberwithin{equation}{subsection}

% Macros
%
\def\lim{\mathop{\mathrm{lim}}\nolimits}
\def\colim{\mathop{\mathrm{colim}}\nolimits}
\def\Spec{\mathop{\mathrm{Spec}}}
\def\Hom{\mathop{\mathrm{Hom}}\nolimits}
\def\Ext{\mathop{\mathrm{Ext}}\nolimits}
\def\SheafHom{\mathop{\mathcal{H}\!\mathit{om}}\nolimits}
\def\SheafExt{\mathop{\mathcal{E}\!\mathit{xt}}\nolimits}
\def\Sch{\mathit{Sch}}
\def\Mor{\mathop{\mathrm{Mor}}\nolimits}
\def\Ob{\mathop{\mathrm{Ob}}\nolimits}
\def\Sh{\mathop{\mathit{Sh}}\nolimits}
\def\NL{\mathop{N\!L}\nolimits}
\def\CH{\mathop{\mathrm{CH}}\nolimits}
\def\proetale{{pro\text{-}\acute{e}tale}}
\def\etale{{\acute{e}tale}}
\def\QCoh{\mathit{QCoh}}
\def\Ker{\mathop{\mathrm{Ker}}}
\def\Im{\mathop{\mathrm{Im}}}
\def\Coker{\mathop{\mathrm{Coker}}}
\def\Coim{\mathop{\mathrm{Coim}}}

% Boxtimes
%
\DeclareMathSymbol{\boxtimes}{\mathbin}{AMSa}{"02}

%
% Macros for moduli stacks/spaces
%
\def\QCohstack{\mathcal{QC}\!\mathit{oh}}
\def\Cohstack{\mathcal{C}\!\mathit{oh}}
\def\Spacesstack{\mathcal{S}\!\mathit{paces}}
\def\Quotfunctor{\mathrm{Quot}}
\def\Hilbfunctor{\mathrm{Hilb}}
\def\Curvesstack{\mathcal{C}\!\mathit{urves}}
\def\Polarizedstack{\mathcal{P}\!\mathit{olarized}}
\def\Complexesstack{\mathcal{C}\!\mathit{omplexes}}
% \Pic is the operator that assigns to X its picard group, usage \Pic(X)
% \Picardstack_{X/B} denotes the Picard stack of X over B
% \Picardfunctor_{X/B} denotes the Picard functor of X over B
\def\Pic{\mathop{\mathrm{Pic}}\nolimits}
\def\Picardstack{\mathcal{P}\!\mathit{ic}}
\def\Picardfunctor{\mathrm{Pic}}
\def\Deformationcategory{\mathcal{D}\!\mathit{ef}}
\usepackage{cleveref}
\crefname{originaldefinition}{original definition}{original definitions}
\Crefname{originaldefinition}{Original definition}{Original definitions}
\makeatletter
\newlabel{definition-fill-in-diagram}{{definition-fill-in-diagram}{source}{}{sourcepointer-1}{}}
\newlabel{definition-fill-in-diagram@cref}{{[originaldefinition][0][]definition-fill-in-diagram}{[1][0][]source}{}{}{}}
\makeatother
\crefname{originallemma}{original lemma}{original lemmas}
\Crefname{originallemma}{Original lemma}{Original lemmas}
\makeatletter
\newlabel{lemma-diagonal-diagonal}{{lemma-diagonal-diagonal}{source}{}{sourcepointer-2}{}}
\newlabel{lemma-diagonal-diagonal@cref}{{[originallemma][0][]lemma-diagonal-diagonal}{[1][0][]source}{}{}{}}
\makeatother
\title{The category of dotted arrows is a groupoid, and a setoid when the diagonal is separated}
\author{The Stacks Project Authors}
\date{stacks-project:stacks-morphisms.tex:local-snapshot-20261012; GNU FDL 1.2 or later}
\begin{document}
\maketitle
\noindent\textit{Editorial note.} All delivered mathematical text is the author's own, copied from the pinned Stacks Project chapter source sources/stacks-algebraic-stacks-morphisms/pinned-original-source.tex; no statement or proof step was rewritten, reordered or removed.\par\medskip
\noindent\textit{Editorial note.} Preamble provenance: the Stacks Project keeps one single shared preamble for every chapter; this card takes it directly from the staged sources/stacks-algebraic-stacks-morphisms/preamble.tex as three verbatim spans (lines 8-12, 14-18 and 154-217), so no derived file was written and nothing was added or redefined. The wrapper class and the base packages are supplied by the assembler.\par\medskip
\noindent\textit{Editorial note.} Bibliography provenance: neither the counted statement nor its proof contains a citation (this lane rescans the whole span before writing the card), so the card carries no bibliography and no entry was reconstructed.\par\medskip
\noindent\textit{Editorial note.} Selection provenance: this card was opened under the file-based control contract (CONTROL.md), with the mechanical pre-card gate run on its statement and proof interval before the card was written.\par\medskip
\noindent\textit{Allowed standard prerequisites (editorial).}\begin{itemize}
\item The situation of the valuative criterion: a 2-commutative square with a chosen 2-morphism and the category of dotted arrows filling it. The statement refers to the defining definition of this chapter (source line 9116). Fixes the objects and morphisms whose category is under consideration.
\item The description of an automorphism of a dotted arrow as a 2-morphism on the first component satisfying two compatibility conditions, and its interpretation as a point of the relative inertia pulled back along the arrow. Computed inside the counted proof (source lines 9122-9145). Turns the two conditions into the statement that the resulting point of the group algebraic space is trivial after restriction to the fraction field.
\item The equivalence between separatedness of the diagonal and the closedness of the identity section of the relative inertia. Invoked at the end of the counted proof (source lines 9146-9147) as a pointer to a lemma of this chapter. Concludes that an automorphism restricting to the identity on the fraction field is itself the identity.
\end{itemize}\medskip
\ifdefined\hypertarget\hypertarget{sourcepointer-1}{}\fi
\noindent\textit{Source reference.} Original definition definition-fill-in-diagram: sources/stacks-algebraic-stacks-morphisms/pinned-original-source.tex, source lines 9116--9116. Navigation pointer only: the counted text refers to the Stacks item definition-fill-in-diagram, which belongs to another chapter of the project and is not part of this card. The pointer carries the target's own label as its number because this chapter's staged snapshot does not contain the other chapter's tag table; its content is not reproduced here.\par\medskip
\ifdefined\hypertarget\hypertarget{sourcepointer-2}{}\fi
\noindent\textit{Source reference.} Original lemma lemma-diagonal-diagonal: sources/stacks-algebraic-stacks-morphisms/pinned-original-source.tex, source lines 3959--3959. Navigation pointer only: the counted text refers to the Stacks item lemma-diagonal-diagonal, which belongs to another chapter of the project and is not part of this card. The pointer carries the target's own label as its number because this chapter's staged snapshot does not contain the other chapter's tag table; its content is not reproduced here.\par\medskip
\begin{lemma}
\label{lemma-cat-dotted-arrows}
In the situation of Definition \ref{definition-fill-in-diagram}
the category of dotted arrows is a groupoid. If $\Delta_f$
is separated, then it is a setoid.
\end{lemma}

\begin{proof}
Since $2$-arrows are invertible it is clear that the category of
dotted arrows is a groupoid. Given a dotted arrow $(a, \alpha, \beta)$
an automorphism of $(a, \alpha, \beta)$ is a $2$-morphism
$\theta : a \to a$ satisfying two conditions. The first condition
$\beta = (\text{id}_f \star \theta) \circ \beta$ signifies that
$\theta$ defines a morphism
$(a, \theta) : \Spec(A) \to \mathcal{I}_{\mathcal{X}/\mathcal{Y}}$.
The second condition
$\alpha = \alpha \circ (\theta \star \text{id}_j)$
implies that the restriction of $(a, \theta)$ to $\Spec(K)$
is the identity. Picture
$$
\xymatrix{
\mathcal{I}_{\mathcal{X}/\mathcal{Y}} \ar[d] & &
\Spec(K) \ar[d]^j \ar[ll]_{(a \circ j, \text{id})} \\
\mathcal{X} & & \Spec(A) \ar[ll]_a \ar[llu]_{(a, \theta)}
}
$$
In other words, if $G \to \Spec(A)$ is the group algebraic space
we get by pulling back the relative inertia by $a$, then
$\theta$ defines a point $\theta \in G(A)$ whose image
in $G(K)$ is trivial. Certainly, if the identity $e : \Spec(A) \to G$
is a closed immersion, then this can happen only if
$\theta$ is the identity.
Looking at Lemma \ref{lemma-diagonal-diagonal}
we obtain the result we want.
\end{proof}

\end{document}
